\documentclass{beamer}
\usepackage{graphicx} % Required for inserting images
\usepackage[T1]{fontenc} % Wichtig für Trennung von Wörtern
\usepackage[utf8]{inputenc} % Erlaubt Umlaute im Code
\usepackage[ngerman]{babel}
\usepackage{bbm}
\usepackage{amsmath}
\usepackage{amssymb}
\usepackage{amsthm}
\usepackage{mathtools}
\usepackage{csquotes}
\usepackage{url}

% Package zum verweisen auf Literatur
\usepackage[style=verbose]{biblatex}
\addbibresource{literatur.bib}

\newcommand{\N}{\mathbbm{N}} % natürliche Zahlen
\newcommand{\Z}{\mathbbm{Z}} % ganze Zahlen
\newcommand{\R}{\mathbb{R}} % reelle Zahlen
\newcommand{\C}{\mathbb{C}} % komplexe Zahlen
\newcommand{\Q}{\mathbb{Q}} % rationale Zahlen


% Einstellungen für Template
\setbeamertemplate{footline}[frame number]
\setbeamertemplate{theorems}[numbered]
\setbeamertemplate{section in toc}[sections numbered]
\setbeamertemplate{subsection in toc}[subsections numbered]
\setbeamerfont{footnote}{size=\tiny}

% Sätze, nummeriert nach Section: Satz 1.1, Satz 1.2
\newtheorem{satz}{Satz}[section]
\newtheorem{defi}[satz]{Definition}
\newtheorem{lem}[satz]{Lemma}
\newtheorem{beispiel}[satz]{Beispiel}
\newtheorem{algorithmus}[satz]{Algorithmus}
\newtheorem{proposition}[satz]{Proposition}
\newtheorem{korollar}[satz]{Korollar}

% Überschrift soll Gliederung folgen
\newif\ifshowtoc
\showtoctrue % Standardmäßig auf "wahr"


\AtBeginSection[]
{
  \begin{frame}{Übersicht}
    \tableofcontents[
      currentsection,
      sectionstyle=show/shaded,
      subsectionstyle=show/show/hide
    ]
  \end{frame}
}

% Automatisierung der Titel der Folien:
\newcommand{\secframe}{\begin{frame}{\insertsectionnumber. \insertsection}}
\newcommand{\subsecframe}{\begin{frame}{\insertsectionnumber.\insertsubsectionnumber \ \insertsubsection}}


\title{Rechnen mit Fehlern}
\author{Ayleen Peschke}
\date{12. Mai 2026}

\begin{document}

\maketitle

\begin{frame}{Gliederung}
    \tableofcontents[hideallsubsections]
\end{frame}

\section{Einführung}
    \subsection{Allgemeines}
      \begin{frame}{Allgemeines}
         \begin{itemize}
             \item Bei numerischen Berechnungsproblemen muss immer mit einem Fehler gerechnet werden
             \item Diese entstehen z.B. durch eine fehlerhafte Eingabe
             \item Es wird zwischen drei Arten von Fehlern unterschieden
         \end{itemize}
      \end{frame} 

    \subsection{Datenfehler}
      \begin{frame}{Datenfehler}
          \begin{itemize}
              \item Der Rechner benötigt zu Beginn Eingabedaten
              \item Die Maschinenzahl ist jedoch begrenzt
              \item Eingabedaten sind häufig Messdaten, diese sind oft gestört
              \item $\Tilde{x}$ beschreibt den gestörten Eingabewert zum korrekten Wert $x$
          \end{itemize}
      \end{frame}

     \subsection{Rundungsfehler}
       \begin{frame}{Rundungsfehler}
           \begin{itemize}
               \item Der Maschinenzahlbereich ist nur begrenzt
               \item Alle Zwischenergebnisse müssen gerundet werden
           \end{itemize}
       \end{frame}

     \subsection{Verfahrensfehler}  
       \begin{frame}{Verfahrensfehler}
           \begin{itemize}
               \item Ein Algorithmus ist eine (theoretisch) unendlich lange Folge
               \item Die Folge konvergiert gegen das gewünschte Ergebnis
               \item Nach endlich vielen Schritten wird das Verfahren abgebrochen
           \end{itemize}
       \end{frame}

       \begin{frame}{Modellfehler}
           \begin{itemize}
               \item Das mathematische Modell beschreibt das Problem nicht immer genau
               \item Dabei entstehen die Modellfehler
           \end{itemize}
       \end{frame}

     \subsection{Darstellungen von Fehlern}
       \begin{frame}{Darstellungen von Fehlern}
           \begin{itemize}
               \item Der absolute Fehler $\Delta x:=|x- \Tilde{x}|$ beschreibt die Abweichung des gestörten Wertes zum eigentlichen Wert. 
               \item Der relative Fehler $\varepsilon:= \frac{\Delta x}{x}$ beschreibt den Anteil des absoluten Fehlers vom exakten Wert. 
           \end{itemize}
       \end{frame}       

\section{Binäre Suche} 
\subsection{Definition}
   \begin{frame}{Binäre Suche}
       \begin{defi}
            Sei $f: \R \to \R$ eine monoton wachsende Funktion und sei $y \in \R$. Existieren die Schranken $L,U \in \R$ mit $f(L) \leq y \leq f(U)$, so lässt sich $f^{-1}(y)$ durch die binäre Suche berechnen: \\
            Sei $m_i= \frac{l_i+u_i}{2}$ der Mittelpunkt des Intervalls $[l_i,u_i]$: \\
            Gilt $f(m_i)< y$, setze $m_i=l_{i+1}$ und fahre mit dem Intervall $[l_{i+1},u_i]$ fort und betrachte $f(m_{i+1})$. \\
            Gilt $f(m_i)>y$, setze $m_i=u_{i+1}$ und fahre mit dem Intervall $[l_i,u_{i+1}]$ fort und betrachte $f(m_{i+1})$. \\ 
            Für die Folge $(m_n)_{n \in \N}$ aller Mittelpunkte gilt $\lim_{n \to \infty} =f^{-1}(y)$.
            Man muss die Funktion $f(x)$ nicht kennen, sondern diese nur auswerten können. Man sagt, $f$ ist durch ein Orakel gegeben. 
       \end{defi}
   \end{frame}

   \begin{frame}{Beispiel}
       \begin{beispiel}
         Wir wählen $y=3$, $f(x)=x^2$ und unsere Schranken mit $L=1$ und $U=2$, da $\sqrt{3} \in [1,2]$: \\ 
        \begin{align*}
         1,5^2=2,25 &\Rightarrow \sqrt{3} \in [1,5,\: 2] \\
         1,75^2=3,0625 &\Rightarrow \sqrt{3} \in [1,5, \: 1,75] \\ 
         1,625^2=2,640625 &\Rightarrow \sqrt{3} \in [1,625, \: 1,75] \\
         1,6875^2=2,84765625 &\Rightarrow \sqrt{3} \in [1,6875, \: 1,75] \\
         1,7134375^2=2,9541015625 &\Rightarrow \sqrt{3} \in [1,71875, \: 1,75] \\
         1,734375^2=3,008056640625 &\Rightarrow \sqrt{3} \in [1,71875, \: 1,734375]
        \end{align*}
     \end{beispiel}
   \end{frame}

   \begin{frame}{Binäre Suche}
       \begin{defi}
           Sei $[l_i,u_i]$ das zu betrachtende Intervall und $m_i= \frac{l_i+u_i}{2}$ dessen Mittelpunkt in der i-ten Iteration. 
           Es gilt $|m_i-\sqrt{3}|\leq \frac{1}{2} \cdot|u_i-l_i|=2^{-i}|u_1-l_1|= 2^{-i}|U-L|$
           Diese Schranke wird a-priori-Schranke gennant. \\
           Im Nachhinein lässt sich eine genauere Schranke angeben, zum Beispiel $|m_i-\sqrt{3}|=\frac{|m_i^2-3|}{m_i+ \sqrt{3}}\leq \frac{|m_i^2-3|}{m_i+l_i}$. Diese Schranke nennt man a-posteriori-Schranke. 
       \end{defi}
  \end{frame}
  
  \subsection{Algorithmus}
  
  \begin{frame}{Algorithmus}
     \begin{algorithmus}
       Eingabe: Ein Orakel zur Berechnung einer monoton wachsenden Funktion $f: \Z \to \R$, $L,U \in \Z$ mit $L\leq U$, $y \in \R$ mit $y \geq f(L)$ \\ 
    Ausgabe: das maximale $n \in \{L,...,U\}$ mit $f(n) \leq y$ 
    \begin{align*}
        &l \leftarrow L \\ 
       &u  \leftarrow  U+1 \\
       &\textbf{while} u>l+1 \textbf{do} \\
       &m \leftarrow \left \lfloor \frac{l+u}{2} \right \rfloor \\
       & \textbf{if} \: f(m) > y \: \textbf{then} \: u \leftarrow m \: \textbf{else} \: l \leftarrow m \\ 
       & \textbf{output} \: l
    \end{align*}
    \label{binäre suche}
    \end{algorithmus}
  \end{frame}

  \begin{frame}{Binäre Suche}
      \begin{satz}
      Der Algorithmus \ref{binäre suche} arbeitet korrekt und terminiert nach $O(log(U-L+2))$ Iterationen.  
     \end{satz}
  \end{frame}

  \begin{frame}{Beweis zu Satz 2.5}
      \begin{proof}
      Es gilt jederzeit $L \leq l \leq u-1 \leq U$ und $f(l) \leq y$, sowie $u>U$ und $f(u)>y$. Also liegt das korrekte Ergebnis immer in $\{l,...,u-1\}$. \\ 
      Außerdem gilt 
      \begin{align*}
        &\max\left\{ \left\lfloor \frac{l+u}{2}\right\rfloor -l-1, u- \left\lfloor \frac{l+u}{2} \right\rfloor -1 \right\} \\
        &\leq \max \left\{ \frac{l+u}{2}-l-1, u- \frac{l+u-1}{2}-1 \right\} \\
        &\leq \frac{u-l-1}{2},
     \end{align*}
     das heißt $u-l-1$ wird mit jeder Iteration mindestens halbiert und die Laufzeit beträgt $O(log(U-L+2))$ Iterationen.
   \end{proof}
  \end{frame}

\section{Fehlerfortpflanzung}
  \begin{frame}{Fehlerfortpflanzung}
      \begin{lem}
        Seien $x,y \in \R\setminus\{0\}$ und $\Tilde{x}, \Tilde{y} \in \R$ Näherungen mit $\varepsilon_x:= \frac{x- \Tilde{x}}{x}$ und $\varepsilon_y:= \frac{y- \Tilde{y}}{y}$. Dann gilt für $\varepsilon_\circ := \frac{x \circ y- (\Tilde{x} \circ \Tilde{y})}{x \circ y}$ mit $\varepsilon \in \{+,-,\cdot,/\}$: 
         \begin{align*}
           &\varepsilon_+= \varepsilon_x \cdot \frac{x}{x+y}+ \varepsilon_y \cdot \frac{y}{x+y}, \\
           &\varepsilon_-= \varepsilon_x \cdot \frac{x}{x-y}- \varepsilon_y \cdot \frac{y}{x-y}, \\
           &\varepsilon_{\cdot}= \varepsilon_x + \varepsilon_y- \varepsilon_x \cdot \varepsilon_y, \\
           &\varepsilon_/ = \frac{\varepsilon_x- \varepsilon_y}{1- \varepsilon_y}
         \end{align*}
       \end{lem}
  \end{frame}

  \begin{frame}{Beweis zu Lemma 3.1}
      \begin{proof}
        Es gilt 
        $\varepsilon_x = \frac{x- \Tilde{x}}{x} \Leftrightarrow \Tilde{x}=x \cdot(1- \varepsilon_x)$ und analog $\Tilde{y}=y \cdot(1-\varepsilon_y)$ 
        \begin{align*}
        \varepsilon_+&= \frac{x+y-(\Tilde{x}+\Tilde{y})}{x+y}= \frac{x+y-(1-\varepsilon_x)\cdot x- (1-\varepsilon_y) \cdot y}{x+y} \\
        &= \varepsilon_x \cdot \frac{x}{x+y}+ \varepsilon_y \cdot \frac{y}{x+y}, \\
         \varepsilon_-: \: erfolgt \: analog
        \end{align*}
        \end{proof}
    \end{frame}
    \begin{frame}{Beweis zu Lemma 3.1}
       \begin{proof}
         \begin{align*} 
          \varepsilon_{\cdot}= \frac{x \cdot y- \Tilde{x}\cdot \Tilde{y}}{x \cdot y} &= \frac{x \cdot y-(1- \varepsilon_x)\cdot x \cdot (1- \varepsilon_y) \cdot y}{x \cdot y} \\ 
          & = 1- (1- \varepsilon_x) \cdot (1- \varepsilon_y)= \varepsilon_x+ \varepsilon_y - \varepsilon_x \cdot \varepsilon_y
        \end{align*}
      \end{proof}
    \end{frame}

    \begin{frame}
      \begin{proof}
        \begin{align*}
        \varepsilon_/= \frac{x/y- \Tilde{x}/ \Tilde{y}}{x/y} &= \frac{x/y- ((1- \varepsilon_x) \cdot x)/((1- \varepsilon_y) \cdot y)}{x/y} \\
        &= 1- \frac{1- \varepsilon_x}{1- \varepsilon_y}= \frac{1- \varepsilon_y- 1+ \varepsilon_x}{1- \varepsilon_y}= \frac{\varepsilon_x-\varepsilon_y}{1-\varepsilon_y}
        \end{align*}
       \end{proof} 
    \end{frame} 
      
  

 \section{Kondition} 
   \subsection{Definition}
     \begin{frame}{Kondition}
         \begin{defi}
           Sei $P \subseteq D \times E$ ein numerisches Berechnungsproblem mit $D,E \subseteq \R$. Sei $d \in D$ eine Instanz von $P$ mit $d \neq 0$. \\
           Wir definieren die \textbf{(relative) Kondition} k(d) von d als
           \begin{align*}
             \limsup_{\varepsilon \to 0}& \left\{ \inf \left\{ \frac{|\frac{e-e'}{e}|}{|\frac{d-d'}{d}|}|e \in E, (d,e) \in P \right\} \right\\
             & \left\vert d'\in D, e'\in E, (d',e') \in P, 0<|d-d'|< \varepsilon \right\}
           \end{align*}
          wobei wir hier die Konvention $\frac{0}{0}=0$ verwenden. \\
          Ist k(d) die Kondition der Instanz $d \in D$, so ist die \textbf{(relative) Kondition} des Problems P 
           \begin{equation*}
             \sup \{k(d) | d \in D, d \neq 0\}.
            \end{equation*}
      \end{defi}
     \end{frame}

    \begin{frame}{Kondition}
        \begin{proposition}
          Sei $f:D \to E$ ein eindeutiges numerisches Berechnungsproblem mit $D,E \subseteq \R$, und sei $d \in D$ eine Instanz mit $d \neq 0$ und $f(d) \neq 0$. Dann ist ihre Kondition gleich 
           \begin{equation*}
             \frac{|d|}{|f(d)|} \cdot \limsup_{\varepsilon \to 0} \left\{ \frac{|f(d')-f(d)|}{|d'-d|} | d'\in D, \: 0<|d'-d|< \varepsilon \right\}
          \end{equation*}
        \end{proposition}
    \end{frame}

    \begin{frame}{Kondition}
        \begin{korollar}
          Sei $f: D \to E$ ein eindeutig numerisches Berechnungsproblem, wobei  $D,E \subseteq \R$ und $f$ differenzierbar sei. Dann ist die Kondition einer Instanz $d \in D$ mit $d \neq 0$ und $f(d) \neq 0$ 
           \begin{equation*}
              k(d) = \frac{|f'(d)| \cdot |d|}{|f(d)|}
           \end{equation*}
        \end{korollar}
    \end{frame}
    
  \subsection{Beispiel}  
     \begin{frame}{Kondition}
         \begin{beispiel}
           \begin{itemize}
             \item Sei $f(x)= \sqrt{x}$. Dann ist die Kondition $k(d)= \frac{\frac{1}{2 \sqrt{x}} \cdot x}{\sqrt{x}}= \frac{1}{2}$. Das heißt die Wurzelfunktion ist gut Konditioniert. 
             \item Sei $f(x)= x+a$ und es gilt $k(x)= |\frac{x}{x+a}|$. Ist $|x+a|<|x|$, so ist die Funktion schlecht konditioniert. 
             \item Sei $f(x)= a \cdot x$. Diese Funktion ist immer gut konditionert, da $k(x)= \frac{|a| \cdot |x|}{|x|}=1$
          \end{itemize}
       \label{kondition}
      \end{beispiel}
     \end{frame}

\section{Fehleranalyse}     
   \begin{frame}{Numerisch stabil}
       Ein Algorithmus heißt numerisch stabil, wenn alle Rechenschritte gut konditioniert sind. Numerisch stabil ist, genau wie gut konditioniert, nicht exakt definiert. 
   \end{frame}

   \subsection{Arten der Fehleranalyse}
     \begin{frame}{Arten der Fehleranalyse}
         \begin{itemize}
           \item \textbf{Vorwärtsanalyse}: Es wird untersucht, wie sich der relative Fehler im Laufe einer Rechnung entwickelt.
           \item \textbf{Rückwärtsanalyse}: Jedes Zwischenergebnis wird als exaktes Ergebnis gewertet und es wird abgeschätzt, wie groß die Eingangsdatenstörung dafür sein müsste. 
        \end{itemize}
     \end{frame}

   \subsection{Beispiel}  
     \begin{frame}{Beispiel}
         \begin{beispiel}
           \begin{itemize}
             \item Vorwärtsanalyse: $x \oplus y = \mathrm{rd}(x+y)= (x+y) \cdot(1+ \varepsilon)$ mit $|\varepsilon| \leq eps(F)$
            \item Rückwärtsanalyse: $x\oplus y= x \cdot (1+ \varepsilon)+y\cdot (1+\varepsilon)=\Tilde{x}+ \Tilde{y}$ mit $|\varepsilon|\leq eps(F)$
         \end{itemize} 
     \vspace{\baselineskip}
    \noindent $eps(F)$ beschreibt hierbei den größtmöglichen Rundungsfehler der Maschine. 
     \end{beispiel}
     \end{frame}

     \begin{frame}{Beispiel}
         Wir betrachten erneut $\sqrt{3}$. \\
         Für $1,734375$ ist die korrekte Eingabe $1,734375^2=3,008056640625=3(1+\varepsilon)$ mit $\varepsilon \leq 0,00269$. \\ 
         Da die Kondition $\frac{1}{2}$ beträgt, kann der Fehler durch $\frac{1}{2}\cdot 0,00269$ beschränkt werden.
     \end{frame}

\section{Newton-Verfahren}
    \begin{frame}{Newton-Verfahren}
        \begin{itemize}
            \item Beschreibt ein weiteres Näherungsverfahren
            \item Eignet sich oft besser als die binäre Suche
            \item Gesucht ist die Nullstelle einer differenzierbaren Funktion $f:\R \to \R$ mit $f' \neq 0$
            \begin{equation*}
              x_{n+1} \leftarrow x_n - \frac{f(x_n)}{f'(x_n)}, \: n=0,1,...
             \end{equation*}
             \item Dabei wird ein „geratener“ Startwert $x_0$ gewählt
        \end{itemize}
    \end{frame}
    
   \subsection{Beispiel}
    \begin{frame}{Beispiel}
        \begin{beispiel}
    Um $\sqrt{3}$ zu berechnen, wähle $f(x)= x^2-a$, wobei $a=3$. Es gilt: 
    \begin{equation*}
        x_{n+1} \leftarrow x_n - \frac{x_n^2-a}{2x_n}= \frac{1}{2} \cdot \left(x_n+ \frac{a}{x_n}\right), n=0,1,...
    \end{equation*}
    Wir starten mit $x_0=1$ und erhalten
    \begin{align*}
        &x_1= \frac{1}{2} \cdot \left( 1+\frac{3}{1}\right)=2 \\
        &x_2= \frac{1}{2} \cdot \left( 2+ \frac{3}{2} \right)= \frac{7}{4}=1,75 \\
        &x_3= 1,732142857... \\
        &x_4= 1,732050810...
    \end{align*}
    \end{beispiel}
    \end{frame}
   \subsection{Konvergenzordnung} 
    \begin{frame}{Konvergenzordnung}
        \begin{defi}
          Sei $(x_n)_{n \in \N}$ eine konvergente Folge reeller Zahlen und $x^*:= \lim_{n \to \infty}x_n$. Existiert eine Konstante $c<1$ und ein $n_0 \in \N$, verringert sich der relative Fehler, so dass 
           \begin{equation*}
             |x_{n+1}-x^*| \leq c \cdot |x_n -x^*|
           \end{equation*}
         für alle $n \geq n_0$ gilt, so hat die Folge (mindestens) Konvergenzordnung 1. \\
         Sei $p>1$. Existiert eine Konstante $c \in \R$, so dass
           \begin{equation*}
             |x_{n+1}-x^*| \leq c \cdot |x_n-x^*|^p
            \end{equation*}
         für alle $n \in \N$ gilt, so hat die Folge (mindestens) Konvergenzordnung p. \\
         Eine Folge konvergiert linear bzw. quadratisch, falls die Konvergenzodnung 1 bzw. 2 ist.
      \end{defi}
    \end{frame}

    \begin{frame}{Konvergenzordnung}
        \begin{satz}
          Sei $a \geq 1$ und $x_0 \geq 1$. Die Folge $(x_n)_{n \in \N}$ ist definiert durch $x_{n+1}= \frac{1}{2}\left(x_n+ \frac{a}{x_n}\right)$ für $n=0,1,...$. Für diese Folge gilt:
           \begin{enumerate}[(\alph*)]
           \item $x_1 \geq x_2 \geq ... \geq \sqrt{a}$
           \item Die Folge $x_n$ konvergiert quadratisch gegen $\lim_{n \to \infty} x_n= \sqrt{a}$. Genauer gilt für alle $n \geq0$:
            \begin{equation*}
            x_{n+1}- \sqrt{a} \leq \frac{1}{2}(x_n-\sqrt{a})^2
            \end{equation*}
       \item Für alle $n \geq 1$ gilt:
          \begin{equation*}
            x_n- \sqrt{a} \leq x_n - \frac{a}{x_n} \leq 2(x_n- \sqrt{a})
        \end{equation*}
       \end{enumerate}
     \end{satz}
    \end{frame}

    \begin{frame}{Beweis Satz 6.3 (a)}
        \begin{proof} 
    \begin{enumerate}[(\alph*)]
        \item Nach dem Satz vom arithmetischen und geometrischen Mittel gilt $\sqrt{xy} \leq \frac{x+y}{2}$. Daraus folgt
        \begin{equation*}
            \sqrt{a}= \sqrt{x_n+\frac{a}{x_n}} \leq \frac{1}{2} \left(x_n + \frac{a}{x_n}\right)=x_{n+1}
        \end{equation*}
        Außerdem gilt für $n\geq 1$:
        \begin{equation*}
            x_{n+1}-x_n= \frac{1}{2} \left(x_n+ \frac{a}{x_n} \right)-x_n= \frac{x_n^2}{2x_n}+ \frac{a}{2x_n}- \frac{2x_n^2}{2x_n}= \frac{1}{2x_n}(a-x_n^2)\leq0,
        \end{equation*}
        da $\sqrt{a} \leq x_n$ äquivalent ist zu $a \leq x_n^2$. Daraus folgt $x_{n+1} \leq x_n$, woraus die Behauptung folgt.
        \end{enumerate}
      \end{proof}
    \end{frame}

    \begin{frame}{Beweis Satz 6.3 (b)}
        \begin{proof}
         \begin{enumerate}[(\alph*)]
         \setcounter{enumi}{1}
              \item Für alle $n \geq 0$ ist:
                \begin{equation*}
                   x_{n+1}-\sqrt{a}= \frac{1}{2} \left(x_n+ \frac{a}{x_n} \right)- \sqrt{a} = \frac{1}{2x_n}(x_n^2+a-2x_n\sqrt{a}).
                \end{equation*}
            Es gilt $\frac{1}{x_n}(x_n-\sqrt{a})\leq 1$ und daraus folgt $x_{n+1}-\sqrt{a} \leq \frac{1}{2}(x_n-\sqrt{a})$ für $n \geq1$. Daraus folgt $\lim_{n \to \infty} x_n = \sqrt{a}$ beweist, da sich der Fehler mit jedem Schritt mindestens halbiert. Aus $x_n \geq 1$ folgt die Behauptung.
            \end{enumerate}
        \end{proof}
    \end{frame}

    \begin{frame}{Beweis Satz 6.3 (c)}
        \begin{proof}
            \begin{enumerate}[(\alph*)]
            \setcounter{enumi}{2}
                \item Für $n \in \N$ gilt:
                  \begin{equation*}
                    x_n-\sqrt{a} \overset{\text{(a)}}{\leq} x_n- \frac{a}{x_n}=  2x_n-2x_{n+1} \overset{\text{(a)}}{\leq} 2(x_n- \sqrt{a}) 
                 \end{equation*}
            \end{enumerate}
        \end{proof}
    \end{frame}

    \begin{frame}{}
        Vielen Dank für eure Aufmrksamkeit! \footcite[Der Vortrag hierzu basiert auf:][]{hougardy2018}
    \end{frame}

    \section*{Literatur}

    \begin{frame}{Literaturverzeichnis}
        \printbibliography
    \end{frame}


\end{document}
